\chapter{Galactic Dynamics}\label{ch:vii}

In the following we're going to discuss about the dynamics of galaxies and all the various ways they can be relaxed through interactions. The interested reader can refer to \cite{2009PhT....62e..56B} for a more thorough analysis and coverage of the subject.

\section{Strong Encounters}

In a large stellar system, gravitational force at any point due to all the other stars is almost constant. A star traces out an orbit in the smooth potential of the whole galaxy, and its orbit will be changed if the star passes very close to another star; in this optic, we then define a strong encounter as one that leads to $\Delta V \sim V$.
\par A strong encounter has happened when the change in potential energy is at least as large as the initial average relative kinetic energy. Let’s take two masses $m$
\begin{equation*}
    G\frac{m^{2}}{r} \approx \frac{1}{2}\mu V^{2} \implies r_{s}\approx \frac{4Gm}{V^{2}}
\end{equation*}
with $V$ the initial speed. What's the timescale associated with a strong encounter?
\par In a time $t$ a strong encounter will occur if any other star intrudes on a cylinder of radius $r_{s}$ being swept out of the orbit. Such cylinder has volume $\mathcal{V} = \pi r_{s}^{2}Vt$, so, provided a stellar density $n$ is given, the typical timescale between encounters is (recall $\tau = 1/\pi r_{s}^{2}Vn$) 
\begin{equation*}
    \tau_{SE} = \frac{V^{3}}{4\pi G^{2}m^{2}n} \approx 4\power{10}{12}\left(\frac{V}{10\unit{km}\unit{s}^{-1}}\right)^{3}\left(\frac{m}{M_{\odot}}\right)^{-2}\left(\frac{n}{1\unit{pc}^{-3}}\right)^{-1}\unit{yr}
    \label{eq:setimescale}
\end{equation*}
This means that stars in the disk of galaxies ($V \sim 30 \unit{km}\unit{s}^{-1}$, $n \sim 0.1 \unit{pc}^{-3}$ near the Sun), never physically collide, and are extremely unlikely to pass close enough to deflect their orbits substantially. On the other hand, in a globular cluster ($V \sim 10 \unit{km}\unit{s}^{-1}$, $n \sim 1000 \unit{pc}^{-3}$ or more), strong encounters will be common (i.e. one or more per star in the lifetime of the cluster).
\par Note that for the purposes of the present discussion, we can consider the stars as \emph{point masses} (since their sizes are small compared to their separations) and the gravitational potential of a galaxy to be the sum of a smooth component (the average over a region containing many stars) and a very deep potential around each individual star. The motion of stars within the galaxy will therefore be determined almost entirely by the smooth part of the force.

\section{Dynamical Friction}

The first description of dynamical friction was that due to Chandrasekhar, back in 1943. It is still an extremely accurate rule of thumb some 80 years later and a good place to start. We're going to follow this strategy:
\begin{enumerate}
    \item Derive the effect of one scattering event;
    \item  Integrate over all impact parameters;
    \item  Integrate over all particles.
\end{enumerate}
We will assume an infinite, homogeneous background distribution of particles. We first consider one star of unitary mass in a static, spherically symmetric gravitational field, which means angular momentum is conserved
\begin{equation*}
    \vec{r}\wedge\der{\vec{r}}{t} = \vec{L} = \text{const.} \to L = r^{2}\dot{\phi}
\end{equation*}
so that the orbit lies in a plane. Using polar coordinates, the force (per unit mass) $F = \ddot{r} - r\dot{\phi}^{2}$ can be expanded as
\begin{equation*}
    F(r) = \frac{L}{r^{2}}\frac{\text{d}}{\text{d}\phi}\left(\frac{L}{r^{2}}\der{r}{\phi}\right)-\frac{L^{2}}{r^{3}}
\end{equation*}
which in terms of a new variable $u = 1/r$ becomes 
\begin{equation*}
    \dder{u}{\phi} + u = \frac{GM}{L^{2}}
\end{equation*}
which is solved by an oscillating motion
\begin{equation}
    u(\phi) = C\cos(\phi-\phi_{0}) +\frac{GM}{L^{2}}
\end{equation}
where $\phi = \phi_{0}$ is clearly the angle when the star is at the pericenter (where $r$ is minimum). Let's now consider two bodies, respectively with mass $M$ and $m$, interacting. In terms of the relative velocity $V$, the velocity of the body of mass $M$ with respect to the center of mass is 
\begin{equation*}
    \Delta \vec{v}_{M} = -\frac{m}{m+M}\Delta\vec{V}
\end{equation*}
Not surprisingly, the problem is reduced to that of a unitary mass in the field of a fixed mass $m+M$. Let us define the following quantities
\begin{equation}
    V_{0} = V(t\to-\infty) \quad L = bV_{0}
\end{equation}
with $b$ the well known impact parameter. Substituting in the equation for $u$ in terms of a new total mass $m+M$
\begin{equation*}
    \frac{1}{r} = C\cos(\phi-\phi_{0}) +\frac{G(M+m)}{b^{2}V_{0}^{2}}
\end{equation*}
which we can differentiate with respect of time 
\begin{equation}
    V = \der{r}{t} = Cr^{2}\dot{\phi}\sin(\phi-\phi_{0}) = CbV_{0}\sin(\phi-\phi_{0})
\end{equation}
where we've used the definition of angular momentum and the fact that its value is constant. Taking the limit $t\to-\infty$ of the original and differentiated equations
\begin{align*}
    -V_{0} &= CbV_{0}\sin(-\phi_{0})\\
    0 &= C\cos(-\phi_{0})+\frac{G(M+m)}{b^{2}V_{0}^{2}}
\end{align*}
The combining of the two yields
\begin{equation*}
    \tan\phi_{0} = -\frac{bV_{0}^{2}}{G(M+m)}
\end{equation*}
If initially the velocity was $\vec{V}_{in} = (V_{0},0)$ and after a deflection by an angle $\theta_{d}$ is $\vec{V}_{out} = (V_{0}\cos\theta_{d},V_{0}\sin\theta_{d})$, then the magnitude of velocity is unchanged, but the vector has changed
\begin{equation*}
    \Delta\vec{V} = \vec{V}_{out}-\vec{V}_{in} = (V_{0}(\cos\theta_{d}-1),V_{0}\sin\theta_{d})
\end{equation*}
therefore
\begin{equation*}
    |\Delta V_{\perp}| = V_{0}\sin\theta_{d} \quad |\Delta V_{\parallel}| = V_{0}(1-\cos\theta_{d})
\end{equation*}
The deflection angle can be easily shown to be $\theta_{d} = 2\phi_{0}-\pi$, meaning
\begin{equation}
      |\Delta V_{\perp}| = \frac{2V_{0}|\tan\phi_{0}|}{1+\tan^{2}\phi_{0}} = \frac{2bV_{0}^{3}}{G(M+m)}\left[1+\frac{b^{2}V_{0}^{4}}{G^{2}(M+m)^{2}}\right]^{-1}
      \label{eq:ortoghonalvelocity_chandra}
\end{equation}
while 
\begin{equation}
    |\Delta V_{\parallel}| = \frac{2V_{0}}{1+\tan^{2}\phi_{0}} = 2V_{0}\left[1+\frac{b^{2}V_{0}^{4}}{G^{2}(M+m)^{2}}\right]^{-1}
    \label{eq:parallelvelocity_chandra}
\end{equation}
Switching back to the lab reference frame, we find for the mass $M$ that 
\begin{align*}
    |\Delta V_{\perp}| &= \frac{2mbV_{0}^{3}}{G(M+m)^{2}}\left[1+\frac{b^{2}V_{0}^{4}}{G^{2}(M+m)^{2}}\right]^{-1}\\
    |\Delta V_{\parallel}| &= \frac{2mV_{0}}{M+m}\left[1+\frac{b^{2}V_{0}^{4}}{G^{2}(M+m)^{2}}\right]^{-1}
\end{align*}
If $M$ is traveling through an homogeneus sea of stars, the average variation of the orthogonal velocity will be zero, since its expression is odd in $b$, whilst the parallel variation will be non-zero, since all of these vectors point in the same direction. The mass $M$ suffers a steady deceleration, known as \emph{dynamical friction}.
\par If the distribution function $f(\vec{v}_{m})$ gives the number of stars encountered with speed between $\vec{v}_{m}$ and $\vec{v}_{m}+\delta\vec{v}_{m}$ per unit volume, then the rate that $M$ encounters these stars is
\begin{equation*}
    2\pi b V_{0}f(\vec{v}_{m})\,\text{d}^{3}v_{m}\diff{b}
\end{equation*}
The variation over time of velocity $\vec{v}_{M}$ is 
\begin{align*}
    \der{\vec{v}_{M}}{t}\Big|_{v_{m}} &= \int_{0}^{b_{\text{max}}} \Delta v_{M\,\parallel}\,2\pi b V_{0}f(\vec{v}_{m})\,\text{d}^{3}v_{m}\diff{b}\\
    &=\vec{V}_{0}f(\vec{v}_{m})\,\text{d}^{3}v_{m}\int_{0}^{b_{\text{max}}}\frac{2mV_{0}}{M+m}\left[1+\frac{b^{2}V_{0}^{4}}{G^{2}(M+m)^{2}}\right]^{-1}2\pi b\diff{b}
\end{align*}
The integral is solved by 
\begin{equation*}
    I = \frac{1}{2}\frac{G^{2}(M+m)^{2}}{V_{0}^{4}}\ln(1+\Lambda^{2}) \quad \Lambda = \frac{b_{\text{max}}V_{0}^{2}}{G(m+M)}
\end{equation*}
Note that the quantity $G(m+M)/V_{0}^{2}$ is the impact parameter corresponding to a 90° deviation, hence we might call it $b_{90}$. The Coulombian parameter $\Lambda$ is then $\Lambda = b_{\text{max}}/b_{90}$. Putting everything back together
\begin{align*}
    \der{\vec{v}_{M}}{t}\Big|_{v_{m}} &= \frac{2mV_{0}}{M+m}\vec{V}_{0}f(\vec{v}_{m})\,\text{d}^{3}v_{m}\frac{1}{2}\frac{G^{2}(M+m)^{2}}{V_{0}^{4}}\ln(1+\Lambda^{2})\cdot 2\pi \\
    &\approx 4\pi G^{2}m(M+m)\ln(\Lambda)\frac{\vec{V}_{0}}{V_{0}^{3}}f(\vec{v}_{m})\,\text{d}^{3}v_{m}
\end{align*}
where we've assumed $\Lambda \gg 1$. We can now recall the definition of $V_{0} = |\vec{V}_{M}-\vec{V}_{m}|$ and integrate over the velocity distribution (assuming $\Lambda$ to be weakly dependent on the velocity)
\begin{equation}
    \der{\vec{v}_{M}}{t} = 4\pi G^{2}m(M+m)\ln(\Lambda)\int f(\vec{v}_{m})\frac{\vec{V}_{m}-\vec{V}_{M}}{|\vec{V}_{m}-\vec{V}_{M}|^{3}}\,\text{d}^{3}v_{m}
    \label{eq:Chandrasekharequation}
\end{equation}
If the velocity distribution $f(\vec{v}_{m})$ is isotropic, the acceleration is given by the total “mass” inside $v_{m}<v_{M}$, hence
\begin{equation*}
    \der{\vec{v}_{M}}{t} = -16\pi G^{2}m(M+m)\ln(\Lambda)\frac{\vec{v}_{M}}{v_{M}^{3}}\int_{0}^{v_{M}} f(\vec{v}_{m})v_{m}^{2}\diff{v_{m}}
\end{equation*}
which is the so-called Chandrasekhar equation\footnote{Note that we're assuming that only stars slower than $V_{M}$ are contributing to slow down $M$.}. We can even assume a Maxwellian distribution for $f$, obtaining
\begin{equation}
    \der{\vec{v}_{M}}{t} = -\frac{4\pi G^{2}\ln\Lambda \rho M}{v_{M}^{3}}\left[\text{Erf}(X)-\frac{2X}{\sqrt{2\pi}}e^{-X^{2}}\right]\vec{v}_{M}
    \label{eq:chandra_max}
\end{equation}
where $X = v_{M}/\sqrt{2\sigma^{2}}$ with $\sigma$ the velocity dispersion and we've assumed $\rho = mN$ under the assumption of $M\gg m$. If the term within squared parentheses tends to 1 (as it does for $v_{M}\gg \sigma$), then 
\begin{equation*}
    \der{\vec{v}_{M}}{t} \sim -\frac{\rho M}{v_{M}^{2}}\hat{v}_{M}
\end{equation*}
If $M$ is large, or $v_{M}$ is small or the density $\rho$ is high, the friction increases!
\par In the opposite regime ($v_{M}\ll \sigma$) we can approximate $f(v_{m})\sim f(0)$ and find
\begin{equation*}
    \der{\vec{v}_{M}}{t} \approx -\rho M \vec{v}_{M}
\end{equation*}
which is the well-known result for Stokes' viscous friction (parachute effect).
\par In short:
\begin{itemize}
    \item $v_{M}\gg \sigma$: This happens when a galaxy or satellite moves through a background whose stars or dark matter have random velocities much smaller than its orbital speed;
    \item $v_{M}\ll \sigma$: A satellite that has already been significantly slowed down, motion in the central regions of a halo or elliptical galaxy--where $\sigma$ is large--or during the late stages of orbital decay. 
\end{itemize}

\subsubsection*{Example: The Decay in the Orbit of a Globular Cluster}

Consider a globular cluster orbiting in a galaxy. Dynamical friction will cause it to lose energy and spiral towards the center. We would like to estimate the $t_{\text{frict}}$ that will take a cluster starting at radius $r_{i}$ to spiral down to the center of the galaxy. Let's assume a flat rotation curve and a dark halo with density (isothermal sphere)
\begin{equation*}
    \rho_{d} = \frac{v_{c}^{2}}{4\pi r^{2}G}
\end{equation*}
For a cluster of mass $M$ at $v = v_{c}$, the force due to the dynamical friction is given by Chandrasekhar's equation (\ref{eq:chandra_max}). The force exterted by friction must match
\begin{equation*}
    \vec{F}_{f} \approx -\frac{\ln\Lambda GM^{2}}{r^{2}}\hat{T}
\end{equation*}
with $\hat{T}$ the unit vector tangent to the orbit. This tangential force will cause the cluster to lose angular momentum at rate
\begin{equation*}
    \der{\vec{L}}{t} = \vec{r}\wedge \vec{F}_{\text{friction}} \approx -\frac{\ln\Lambda GM^{2}}{r}\hat{z}
\end{equation*}
from which 
\begin{equation*}
    r\der{r}{t}\approx -\frac{\ln\Lambda GM}{v_{c}}
\end{equation*}
Integrating from the initial radius $r_{i}$ and the final radius 0, we find $t_{\text{frict}}$
\begin{equation}
    t_{\text{frict}} \approx \frac{r_{i}^{2}v_{c}}{GM\ln\Lambda}\approx \frac{100}{\ln\Lambda}\left(\frac{r_{i}}{2\unit{kpc}}\right)^{2}\left(\frac{v_{c}}{250\unit{km}\unit{s}^{-1}}\right)\left(\frac{10^{6}M_{\odot}}{M}\right)\unit{Gyr}
    \label{eq:friction_ts}
\end{equation}
\par A plausible value for $\Lambda$ would give $\ln\Lambda \sim 10$. In M31 all GCs starting within $3\unit{kpc}$ radius should have sunk to the center by now, and indeed it is possible that the bright nucleus of M31 might be composed by 20-30 massive GCs. Moreover, if we consider two galaxies of the same mass colliding ($M\sim 10^{10}M_{\odot}$, $r_{i}\sim 20\unit{kpc}$, $v = v_{c}$), we'd get an infall time of about $\sim 0.2\unit{Gyr}$, which is roughly equivalent to a single orbital period. Clearly, massive galaxies entering each other's halos experience strong dynamical friction.
\par On a concluding note, dynamical friction may be derived heuristically considering the wake left by the massive object as it moves in the sea of stars. If you're interested, you can take a look at the course slides.

\section{Two-Body Relaxation}

"Collisional or not collisional, that is the question," that's what Prince Hamlet would probably say if he'd decided to dedicate his life to galaxy astrophysics rather than to exact revenge against his uncle Claudius.
\par \emph{Collisional} systems in astrophysics are those in which the mutual gravitational interactions between individual bodies (often through two-body encounters) significantly affect the system's long-term dynamics. These systems are opposed to \emph{collisionless} systems, where the smooth gravitational potential dominates over individual encounters.
\par When can we claim with certainty that a system is, say, collisional rather than collisionless? It depends on the \emph{two-body relaxation timescale.}
\par Earlier we've found that 
\begin{align*}
    |\Delta V_{\perp}| &= \frac{2mbV_{0}^{3}}{G(M+m)^{2}}\left[1+\frac{b^{2}V_{0}^{4}}{G^{2}(M+m)^{2}}\right]^{-1}\\
    |\Delta V_{\parallel}| &= \frac{2mV_{0}}{M+m}\left[1+\frac{b^{2}V_{0}^{4}}{G^{2}(M+m)^{2}}\right]^{-1}
\end{align*}
Averaging over many interactions, this would imply that the mean effect of parallel encounters will always reduce the velocity of the heavier object (dynamical friction), while the mean effect of perpendicular encounters will average to zero. Even though there is no net shift in the perpendicular direction, there will still be a spread in the $\Delta V_{\perp}$ values, which can be measured in terms of the standard deviation or variance. This spread quantifies the "energy" associated with the perpendicular components of the satellites' velocities
\begin{equation*}
    \sigma_{\perp}^{2} = \langle v_{\perp}^{2}\rangle - \langle v_{\perp}\rangle^{2} = \langle v_{\perp}^{2}\rangle
\end{equation*}
The average speed will be 0 at any time, but there will be higher and lower speeds depending on the deflection angle. The result is that the velocity dispersion increases (the distribution of velocities broadens). Note that diffusion is more effective in the perpendicular direction. In fact, for small deflection angles
\begin{equation*}
    |\Delta V_{\perp}| = V_{0}\sin\theta_{d} \approx V_{0}\theta_{d} \quad |\Delta V_{\parallel}| = V_{0}(2-\cos\theta_{d})\approx \frac{1}{2}V_{0}\theta_{d}^{2}
\end{equation*}
Hence for small angles $|\Delta V_{\perp}| \gg |\Delta V_{\parallel}|$. A system is "relaxed" when during its lifetime a typical star loses memory of its initial kinetic energy.
\input{chapters/impactparamter.tex}
\par Given this setup, depicted in Fig.\ref{fig:star_impact}, we're now going to show that stars in clusters can reach equilibrium through mutual interactions in a process named \emph{two-body relaxation}, which is essentially analogous to thermalization. Therefore, a very important timescale in collisional dynamics in clusters is the time for a star to completely lose memory of its initial velocity by means of (weak) gravitational encounters.
\par The gravitational force is 
\begin{equation*}
    \vec{F} = -\frac{Gm^{2}}{r^{3}}\vec{r} = \vec{F}_{\parallel}+\vec{F}_{\perp}
\end{equation*}
where the orthogonal component is
\begin{equation*}
    F_{\perp} = -\frac{Gm^{2}}{r^{2}}\frac{b}{r} = -G\frac{m^{2}}{b^{2}}\left[1+\left(\frac{vt}{b}\right)^{2}\right]^{-3/2} = -m\dot{v}_{\perp}
\end{equation*}
since $r^{2} = b^{2}+(vt)^{2}$. The change in velocity integrated over one entire encounter\footnote{Note that we're (a) assuming a straight line trajectory (each interaction has a small effect on the orbit, hence $\delta v_{\parallel}\approx 0$), (b) the secondary star is assumed stationary, and (c) all stars have the same mass $m$.} is 
\begin{equation*}
    \delta v_{\perp} = \int_{-\infty}^{\infty}\dot{v}_{\perp}\,\text{d}t = \frac{Gm}{b^{2}}\int_{-\infty}^{\infty}\left[1+\left(\frac{vt}{b}\right)^{2}\right]^{-3/2}\,\text{d}t = \frac{2Gm}{bv}
\end{equation*}
\par Now taking into account all stars in the system, the surface density of stars in an idealized cluster is $N/\pi b_{\text{max}}^{2}$, and the number of interactions per unit element is 
\begin{equation*}
    \delta n = \frac{N}{\pi b_{\text{max}}^{2}}\text{d}(\pi b^{2}) = \frac{2bN}{b_{\text{max}}^{2}}\text{d}b
\end{equation*}
Defining $\delta v^{2}_{\text{tot}} = \int \delta v^{2}_{\perp}\,\delta n$, where the integral is performed over all possible impact parameters, we get
\begin{equation*}
    \delta v^{2}_{\text{tot}} = 8N\left(\frac{Gm}{b_{\text{max}}v}\right)^{2}\log\left(\frac{b_{\text{max}}}{b_{\text{min}}}\right)
\end{equation*}
The integration limit $b_{\text{max}}$ is of the order of the size of the system, while $b_{\text{min}}$ corresponds to the smallest $b$ to avoid stellar collisions, which is equal to 
\begin{equation*}
    b_{\text{min}} = \frac{2Gm}{v^{2}} = b_{90}
\end{equation*}
thus leading to 
\begin{equation*}
     \delta v^{2}_{\text{tot}} = 8N\left(\frac{Gm}{b_{\text{max}}v}\right)^{2}\log\left(\frac{b_{\text{max}}}{b_{90}}\right)
\end{equation*}
The typical speed of a star in a virialized system\footnote{We're referring to stars that obey the virial theorem and the equation that follows is merely a restating of the more general theorem.} is given by
\begin{equation*}
    Nmv^{2} = \frac{G(Nm)^{2}}{b_{\text{max}}}
\end{equation*}
Replacing for $v$ in the equation for $\delta v_{\text{tot}}$ we get 
\begin{equation*}
    \delta v^{2}_{\text{tot}} = \frac{8}{N}v^{2}\log\left(\frac{N}{2}\right)
\end{equation*}
The number of crossings for which $\delta v^{2}_{\text{tot}}/v^{2} \approx 1$ (that is the number of crossings after which the star has changed its initial velocity completely and lost memory of its initial conditions) is given by
\begin{equation*}
    n_{\text{cross}}\approx\frac{N}{8}\frac{1}{\log(N/2)}
\end{equation*}
The time needed to cross the system is 
\begin{equation*}
    \tau_{\text{cross}} = \frac{b_{\text{max}}}{v} = \sqrt{\frac{b_{\text{max}}^{3}}{GNm}} \propto \frac{1}{\sqrt{G\rho}}\propto \tau_{ff}
\end{equation*}
while the time necessary for stars in a system to completely lose the memory of their initial velocity, called the \emph{relaxation time}, is
\begin{equation}
    \tau_{rlx} = n_{\text{cross}}\tau_{\text{cross}} = \frac{N}{8}\frac{1}{\log(N/2)}\frac{b_{\text{max}}}{v}
    \label{eq:relaxation_time}
\end{equation}
For globular clusters, dense young star clusters, nuclear star clusters (far from SMBH influence radius) the typical relaxation time is $\approx 10^{7}-10^{10}\,\text{yr}$, thus making them collisional systems. On the other hand, galaxy field/disks have a relaxation time that is rather larger than the Hubble timescale, and are thus collisionless.

\subsection{Effects of two-body relaxation: Mass Segregation}

\par Because of two-body relaxation there is a transfer of kinetic energy (equipartition). If all stars have the same mass, this implies that at equilibrium the velocity distribution tends toward an approximately Maxwellian distribution. Given that $\bar{v}^{2} = 3kT/m$, this implies $\bar{v}^{2} \approx \text{const.}$ For different masses, a Maxwellian distribution is reached for each mass type, but the same temperature will be shared among all masses. In this case, $\bar{v}^{2}_{i} = 3kT/m_{i}$, meaning $m_{i}\bar{v}^{2}_{i}\approx \text{const.}$ If $m_{i}>m_{j}$, then $\bar{v}^{2}_{i} < \bar{v}^{2}_{j}$, and so massive stars slow down, light stars move to higher velocities. We achieve what is commonly known as \emph{mass segregation}. A nice criterion for the determination of mass segregation was given by Spitzer, but we won't discuss it here. Just know that usually mass segregation occurs after about $\sim 1 \tau_{\text{rlx}}$ and it's actually observed in objects like the Pleiades star cluster, see Fig.\ref{fgr:pleiades_massseg}.
\begin{figure}[h!]
    \centering 
    \includegraphics[width=0.7\textwidth]{img/pleiades_massseg.png}
    \caption{Mass segregation in the Pleiades star cluster. Note how massive stars are concentrated at the center, while lower mass stars populate the outskirts.}
    \label{fgr:pleiades_massseg}
\end{figure}
\par Note that this applies if the age of the system is larger than the relaxation timescale.

\subsection{Effects of two-body relaxation: Evaporation}

Evaporation of a cluster is yet another process able to disrupt a cluster. Describing the star cluster as a self-gravitating system, we can define an escape velocity for any given distance from its center 
\begin{equation*}
    v^{2}_{e}(r) = -2\phi(r)
\end{equation*}
The mean square escape velocity is then obtained by averaging over the (spherically symmetric) mass density distribution $\rho(r)$
\begin{equation*}
    \langle v^{2}_{e}\rangle = \frac{\int \rho(r)v^{2}_{e}(r)\,\text{d}r}{\int \rho(r)\,\text{d}r} = -\frac{2}{M}\int \rho(r)\phi(r)\,\text{d}r = -\frac{4E_{g}}{M}
\end{equation*}
where $E_{g}$ is the potential self-energy and $M$ is the total mass. Using the virial theorem once more leads to 
\begin{equation*}
    \langle v^{2}_{e}\rangle = 4\langle v^{2}\rangle
\end{equation*}
Therefore, stars with velocities exceeding twice the root-mean-square (RMS) velocity of the distribution are unbound. For a typical Maxwellian velocity distribution, this amounts to a fraction $\epsilon = 7.4\cdot 10^{-3}$ of all stars.
\par Roughly, evaporation removes $\text{d}N = -\epsilon N$ stars on a relaxation timescale, so the system gets slightly hotter and contracts. The velocity distribution adjusts to another Maxwellian distribution, so that in every relaxation time $\epsilon N$ stars are removed by evaporation
\begin{equation}
    \frac{\text{d}N}{\text{d}t} = -\frac{\epsilon N}{t_{rlx}} := -\frac{N}{t_{\text{evap}}}
    \label{eq:evaporation_ts}
\end{equation}
Note that the evaporation timescale is much greater than the relaxation timescale. Consider again the Pleiades star cluster. The relaxation timescale for the Pleiades is about $160\unit{Myr}$ while the evaporation timescale is roughly $\sim 20\unit{Gyr}$. Compare it now with the age of the cluster, which is estimated to be about $125 \unit{Myr}$. Clearly, we can expect it to be decently mass-segregated, but the cluster has not yet lost much mass due to evaporation.
\par Now, evaporation is not only mechanism for clusters to lose mass. For example, external gravitational perturbations to star clusters (from GMCs, disk of Milky Way, etc) also pump energy into stellar orbits (tidal heating) or directly strip them (tidal stripping), hastening the mass loss. Star clusters with mass like Pleiades that formed early in galaxy's history with age $\sim 10\unit{Gyr}$ have already dissolved through evaporation and/or tidal heating.
\par Lower mass clusters dissolve even more rapidly than the Pleiades. The only star clusters surviving from the formation time of the galaxy are massive star clusters ($N > 10^{5}$ stars), i.e. globular clusters.
\par Another effect of two-body relaxation (which is not relevant in galaxies) is \emph{core collapse}: Relaxation is more effective in the more dense parts of a system (core). Since the core is more collisional, more evaporation is due, and hence more energy is removed from the system. If you remove energy (heat), stars fall deeper in the gravitational well (so the system contracts getting denser), they therefore speed up, and that part of the system gets hotter, thus leading to more evaporation.
\par Mass loss due to evaporation then leads to collapse, a runaway process called \emph{gravothermal instability}. If the system contracts, it becomes denser and the two-body encounter rate increases along with the evaporation rate. As a consequence, the core of the cluster loses energy (the kinetic energy of the evaporated stars) to the halo, and $\text{d}E_{\text{core}}<0$. Since for any bound, finite system in which the dominant force is gravity we have
\begin{equation*}
    C := \frac{\text{d}E}{\text{d}T} < 0
\end{equation*}
the temperature of the core must increase. Therefore, stars exchange more energy and become dynamically hotter, and faster stars tend to evaporate at a higher rate. This runaway process leads to the unphysical situation of star clusters with infinite core density.
\par To avoid this catastrophic scenario, we consider the possibility that the core collapse is reversed by an external source injecting kinetic energy into the core ($\text{d}E_{\text{core}} > 0)$, thus cooling it ($\text{d}T_{\text{core}} < 0$) until the temperature gradient declines to zero, thus halting the collapse. There are multiple ways this can happen: Mass loss by stellar winds and/or SNs; formation of binaries; three-body encounters between single stars and binaries extracting kinetic energy from the internal energy of the binary system and so on.
\par It turns out that core collapse is eventually halted by binary stars. In most clusters these will be primordial, but even if no primordial binaries are present, binaries are formed by encounters of three stars in dense cores. The so-called \emph{hard binaries}, i.e. those with binding energies larger than the average kinetic energy of stars in the cluster, tend to lose energy in such encounters. They become ever more bound and are more likely to provide energy to the next star that passes by. A single hard binary might even have a binding energy equal to that of the cluster as a whole!
\par This is all nice and dandy, but it's not actually that helpful to explain why most galaxy we observe are "relaxed." Clearly it can't be because of the two-body relaxation.
\begin{quote}
    \emph{Galaxy changes are by revolution rather than by evolution.}
\end{quote}
\par Consider a star that is at rest at the center of a collapsing spherical protogalaxy. As the protogalaxy collapses, the potential well at its center becomes deeper. On the other hand, the velocity of the central star remains zero. Therefore the energy of this star decreases.
\par Other stars in the collapsing protogalaxy will gain energy. For example, consider a star that is initially located outside the half-mass radius of the system and is moving slowly radially outward. This star will be slow to respond to the overall collapse of the system, and by the time it is falling rapidly to the center, the system as a whole will be approaching its most compact configuration. Hence the star will acquire a lot of kinetic energy as it falls into the deep potential well at the center of the system.
\par Later, by the time the star has passed close to the center and is on its way out again, the system will have re-expanded significantly and the potential well out of which it has to climb will be less deep than that into which it fell. Consequently, it will reach the potential at which it originally started with more kinetic energy than it had originally. Therefore the energy of this star increases.

\section{Violent relaxation}

During the collapse of large cold systems, the gravitational potential changes rapidly. As a consequence, stars gain and lose energy, which broadens the energy distribution. Energy is then redistributed via collective interactions: This acts like a relaxation process.
\par Note that the total energy remains constant: This is a non-dissipational process. Energy is not radiated away, as with dissipational (gaseous) collapse. If the total energy is initially zero (e.g. diffuse system at rest), then following collapse some stars will be strongly bound, but some must also have been ejected.
\par Since the acceleration of a particle in the mean gravitational field is independent of its mass, the equilibrium configuration is also independent of mass. No segregation by mass will occur. Moreover, violent relaxation is fast: The timescale is the dynamical timescale (\ref{eq:ff_timescale}) $\propto 1/\sqrt{G\rho}$.
\par Numerical simulations, such as that in Fig.\ref{fgr:violentrelaxation}, show that the density distribution becomes smooth before scattering is complete, hence relaxation ceases and is incomplete (only partially phase-mixed in velocity/energy space).
\begin{figure}[h!]
    \centering
    \includegraphics[width=0.6\textwidth]{img/violentrelaxation.png}
    \caption{Four stages of a collapse simulation. Credits: Unknown.}
    \label{fgr:violentrelaxation}
\end{figure}
\par Why does relaxation cease prematurely? Violent relaxation depends on the time-dependent potential to shuffle energies. When that time dependence becomes small (the bulk of the mass is in a more stable configuration), there is no further “violent” driver to continue randomizing orbits. This effectively freezes the system in a partially relaxed state.

\subsection{Scalar Virial Theorem}

Consider autogravitating stellar system, onto which an external force $F_{\text{ext}}$ is applied; this might represent, for example, the gravitational pull of a galaxy on a star cluster within it.
\par The force seen by star $\alpha$ is 
\begin{equation*}
    \der{(m_{\alpha}\vec{v}_{\alpha})}{t} = -\sum_{\beta} \frac{Gm_{\alpha}m_{\beta}(\vec{r}_{\alpha}-\vec{r}_{\alpha})}{|\vec{r}_{\alpha}-\vec{r}_{\alpha}|^{3}}+\vec{F}_{\text{ext}}
\end{equation*}
We can project on the position of the star $\alpha$ by taking the scalar product
\begin{equation*}
    \sum_{\alpha}\der{(m_{\alpha}\vec{v}_{\alpha})}{t}\cdot \vec{r}_{\alpha} = - \sum_{\alpha}\sum_{\beta}\frac{Gm_{\alpha}m_{\beta}(\vec{r}_{\alpha}-\vec{r}_{\alpha})}{|\vec{r}_{\alpha}-\vec{r}_{\alpha}|^{3}}\cdot \vec{r}_{\alpha}+\sum_{\alpha}\vec{F}_{\text{ext}}\cdot \vec{r}_{\alpha}
\end{equation*}
I'm definitely skipping all the intermediate passages because it's pure algebra and estimating time averages. What you come out with is
\begin{equation*}
    \frac{1}{2}\dder{I}{t} - 2K = \Omega +\sum_{\alpha}\vec{F}_{\text{ext}}\cdot \vec{r}_{\alpha}
\end{equation*}
As long as the system is stationary (or stationary on average), the derivative of the moment of inertia is finite and goes to zero when averaging\footnote{Physically speaking, if the system were collapsing $I$ would decrease systematically, if it were expanding, $I$ would increase.}. As a consequence
\begin{equation}
    2\langle K \rangle + \langle \Omega \rangle = -\sum_{\alpha}\langle \vec{F}_{\text{ext}}\cdot \vec{r}_{\alpha} \rangle \to 0 \; (\text{if the system is isolated})
    \label{eq:virial_theorem}
\end{equation}
The conditions to use this equation are simple: The system must be gravitationally bound, the system must be in steady state (so the orbital timescale $\ll$ the evolution timescale), quantities must be time averaged (or many objects sampled with random orbital phase) and the system must be isolated (or at least embedded in a slowly varying potential).

\section{Galaxy Encounters}

During a tidal encounter, the orbits of many stars are significantly affected. However, some orbits are not greatly affected, those for which $t_{\text{orbit}}\ll t_\text{enc}$. As the tidal field slowly changes, the orbit responds approximately reversibly. This type of response is called \emph{adiabatic}. Stars on rapid orbits near galaxy centers are not greatly affected by tidal encounters (unless, of course, the encounter proceeds to become a merger).
\par The opposite extreme occurs when $t_{\text{orbit}}\gg t_\text{enc}$, this occurs when $v_{\text{int}}\ll v_{\text{enc}}$. In this case, stars don't move much during the encounter, hence there's no change in potential energy. However, they do feel an impulse, and thus the kinetic energy will have to change accordingly. If this happens, we say that the system has experienced a \emph{tidal shock}.

\subsection{High Speed Encounter}

In case the encounter velocity is much larger than the internal velocity dispersion of the perturbed system, the change in the internal energy can be obtained analytically using the \emph{impulse approximation}. Consider the encounter between S (the source) and P (perturber). Under the impulse approximation we may consider a particle q in S to be stationary (with respect to the center of S) during the encounter; q only experiences a velocity change $\Delta v$, but its potential energy remains unchanged
\begin{equation*}
    \Delta K_{q} = \frac{1}{2}(\vec{v}+\Delta \vec{v})^{2} - \frac{1}{2}\vec{v}^{2} = \vec{v}\Delta\vec{v} + \frac{1}{2}\Delta\vec{v}^{2}
\end{equation*}
We are interested in computing $\Delta K_{S}$, which is the integral of $\Delta K_{q}$, over the entire system S. Because of symmetry, the term $\vec{v}\Delta\vec{v}$ will equal zero
\begin{equation*}
    \Delta K_{S} = \frac{1}{2}\int \Delta\vec{v}^{2}\rho(r)\,\text{d}^{3}r
\end{equation*}
We set our coordinates so that the velocity of the perturber lies along the horizontal $z$ axis, and the perturber’s orbit lies in the $yz$ plane. In the limit of very high $V$, the perturber moves on the straight line with impact parameter $b$, so that its trajectory is simply $R(t)\approx (0, b, Vt)$.
\par The potential due to the perturber P at the point q is
\begin{equation*}
    \Phi_{P}  = -\frac{GM_{P}}{|\vec{r}-\vec{R}|}
\end{equation*}
where $|\vec{r}-\vec{R}|$ can be estimated from the cosine theorem.
\par Let's expand $\Phi_{p}$ in terms of a small $\epsilon = r/R$. For this purpose, we also define an auxiliary variable $u = -2\epsilon\cos\phi +\epsilon^{2}$. In these terms, the potential becomes
\begin{equation*}
    -\frac{GM_{P}}{|\vec{r}-\vec{R}|} \approx \left[ -\frac{GM_{P}}{R}-\frac{GM_{p}r}{R^{2}}\cos\phi-\frac{GM_{P}r^{2}}{R^{3}}\left(\frac{3\cos^{2}\phi -1}{2}\right)\right]
\end{equation*}
The third term describes the tidal force, which is that we're interested into. Taking the gradient of the potential, and dropping the second (constant acceleration) term, yields the tidal force per unit mass. Integrating the force over time yields the cumulative change in velocity with respect to the center of S. After some more algebra (see \cite{2009PhT....62e..56B} for more) one finds that
\begin{equation*}
    \Delta\vec{v} = \frac{2GM_{p}}{Vb^{2}}(-x,y,0)
\end{equation*}
Finally, integrating over the entire system S and assuming spherical symmetry\footnote{In particular, we can make use of the following identity: $\langle x^{2}+y^{2}\rangle = \langle 2r^{2}/3\rangle$.}
\begin{equation}
    \Delta K_{S} = \frac{4G^{2}M_{P}^{2}M_{S}\langle r^{2}\rangle}{3V^{2}b^{4}}
    \label{eq:tidalshock}
\end{equation}
\par Closer encounters have a much greater impact, but you should recall that our derivation is only valid for a distant tide.
\par How does the system respond (relax) after experiencing the tidal shock? We can use the virial theorem, which will apply to the original ("o" subscript) and the final relaxed system ("f" subscript) but not to the in-between ("i" subscript). Hence, since $E_{o} = -K_{o}$ and $E_{f} = -K_{f}$
\begin{equation*}
    K_{i} = K_{o}+\Delta K \;\text{and}\; E_{i} = E_{o}+\Delta K = -K_{o}+\Delta K
\end{equation*}
while post the final relaxation
\begin{equation*}
    E_{f} = E_{i} \implies -K_{f} = -K_{o}+\Delta K \implies K_{f} = K_{o}-\Delta K
\end{equation*}
What we gather from this is that from the original to the final configuration, the system has indeed cooled, by an amount $2\Delta K$. Since the shock heats the original system by $\Delta K$, then during relaxation (i to f) the system cools by $-2\Delta K$, and in fact $K_{i}=K_{o}+\Delta K$. As a consequence of the virial theorem the gravitational energy is now "less" negative (i.e., the system expands).
\par Before virialization, since the stars receive energy, some may become unbound ($E > 0$). These stars evaporate and are lost from the system. If there are repeated tidal shocks, a star cluster (for example) may be disrupted and disintegrate. On a conclusive note, ring galaxies are said to be formed from the tidal shocks we've just discussed.

\section{Tides}

In astronomy, a simplified problem of importance is the restricted three-body problem, in which two massive particles ($M_{1}$, $M_{2}$) exert gravitational pull on a low-mass test particle ($m$) that doesn’t influence the others.
\par As a further simplification, we'll assume these orbits to be circular. This is usually a good approximation for binary systems, since tidal effects tend to circularize originally eccentric orbits on timescales short compared to the time over which mass transfer occurs.
\par For the purposes of this section, we're going to assume that the two stars of the binary can be assumed to be point-like sources of potential with separation $a = |\vec{r}_{1}-\vec{r}_{2}|$.
\par That being said, the binary period\footnote{Here we're using the angular frequency rather than the period. It's the same up to some power of $2\pi$.} is inferred from Kepler's third law
\begin{equation}
    \Omega^{2} = \frac{G(m_{1}+m_{2})}{a^{3}}
    \label{eq:kepler_3}
\end{equation}
An aesthetically pleasing depiction of the system at hand is shown in Fig.\ref{fig:binaries}.
\begin{figure}[h!]
    \includegraphics[width=0.8\textwidth]{img/binary_pleasing.png}
    \caption{A binary system of two stars orbiting around their common center of mass. Credits: \cite{Frank_King_Raine_2002}.}
    \label{fig:binaries}
\end{figure}
\par In the corotating frame, the potential generated from the system is 
\begin{equation*}
    \phi_{\text{eff}}(\vec{r}\,) = -\frac{Gm_{1}}{|\vec{r}-\vec{r}_{1}|} - \frac{Gm_{2}}{|\vec{r}-\vec{r}_{2}|} -\frac{1}{2}(\vec{\Omega}\wedge\vec{r}\,)^{2}
\end{equation*}
where $\vec{r_{1}}$, $\vec{r_{2}}$ are the position vectors of the centres of the two stars. This is sometimes known as \emph{Roche potential}.
\par We’re looking for equilibrium points of the system, that means we’re looking for properties for which particle $m$ is at rest in the rotating system, so $\ddot{r} = 0$ and $\dot{r} = 0$. Without the Coriolis forces, the equation of motion is just $\ddot{\vec{r}} = -\nabla\phi_{\text{eff}}$; the equilibrium positions will then satisfy $\nabla\phi_{\text{eff}}=0$. Analitical solution is generally not possible, since a 5th order polynomial is involved.
\par It is worth remembering, however, the presence of 5 points at which the gradient of the Roche potential vanishes, which are called Lagrange points $L_{i}$. The three points $L_{1}\to L_{3}$ that lie on the line that joins the two stars are all unstable, while $L_{4}$ and $L_{5}$ are stable. As the ratio between the two masses $q = M_{2}/M_{1}$ becomes small, the Lagrange points $L_{1}$ and $L_{2}$ shrink to roughly equal distances on either side of the secondary (the less massive star).
\begin{figure}[h!]
    \centering 
    \includegraphics[width=0.7\textwidth]{img/roche_sol.png}
    \caption{Equipotential surfaces for the Roche potential $\phi_{\text{eff}}$ for decreasing values of $q$.}
\end{figure}
\par In the limit $q \ll 1$ it's possible to derive the position of the Lagrange point $L_{1}$. We know that for this in particular $y = 0$ and $x_{2}<x< x_{1}$ (the Lagrange points all lie in the $xy$ plane because of angular momentum conservation). Hence
\begin{equation*}
    \tilde{\phi}_{\text{eff}} = -\frac{1}{1+q}\frac{1}{\sqrt{\left[\tilde{x}-\frac{q}{1+q}\right]^{2}+\tilde{y}^{2}}}-\frac{q}{1+q}\frac{1}{\sqrt{\left[\tilde{x}+\frac{1}{1+q}\right]^{2}+\tilde{y}^{2}}}-\frac{1}{2}(\tilde{x}^{2}+\tilde{y}^{2})
\end{equation*}
where we've defined 
\begin{equation*}
    \tilde{x} = \frac{x}{a} \quad \tilde{y} = \frac{y}{a} \quad \tilde{\phi}_{\text{eff}} = \phi_{\text{eff}}\frac{a}{G(M_{1}+M_{2})}
\end{equation*}
and also 
\begin{equation*}
    \tilde{x}_{1} = \left(\frac{q}{q+1}\right) \quad \tilde{x}_{2} = -\left(\frac{1}{q+1}\right)  \implies \tilde{x}_{1}-\tilde{x}_{2} = 1
\end{equation*}
For $q \ll 1$ the gradient of the adimensional potential becomes
\begin{equation*}
    \parder{\tilde{\phi}_{\text{eff}}}{\tilde{x}} = -\frac{1}{(q+1)(\tilde{x}_{1}-\tilde{x})^{2}}+\frac{q}{(q+1)(\tilde{x}_{2}-\tilde{x})^{2}}-\tilde{x} = 0
\end{equation*}
We can define a new variable $\tilde{r} = \tilde{x}-\tilde{x}_{2}$ that is the dimensionless distance from $M_{2}$, thus obtaining 
\begin{equation*}
    -\frac{1}{(1-\tilde{r})^{2}}+\frac{q}{\tilde{r}^{2}}+1-(1+q)\tilde{r} = 0
\end{equation*}
If $q \ll 1$, then $\tilde{r}\ll 1$, and so we can perform some Taylor expansions. In conclusion, we define Hill (or Jacobi) radius as 
\begin{equation}
    \boxed{
    \frac{r_{\text{H}}}{a} = \left(\frac{q}{3}\right)^{1/3}
    }
    \label{eq:hillradius}
\end{equation}
which is telling us that the (average) densities within the Hill radius are in the ratio 
\begin{equation*}
    \rho_{2} = 3\rho_{1}
\end{equation*}
Given that the orbital period is of the order of $1/\sqrt{\rho}$, this means that 
\begin{equation*}
    T_{2} < T_{1}
\end{equation*}
where $T_{2}$ is the period of a star orbiting $M_{2}$, while $T_{1}$ is the period of $M_{2}$ orbiting $M_{1}$.
\par The period of a star orbiting the satellite (secondary) at distance $r_{\text{H}}$ is similar to the satellite orbital period around the main galaxy. The satellite can retain those stars close enough to circle it in less time than it takes to complete its own orbit about the main galaxy.

\section{Ram Pressure Stripping}

A galaxy passing through the ICM will inevitably feel an external pressure. Consider a disk galaxy of radius $R_{d}$ moving through an ICM of density $\rho_{\text{ICM}}$ with relative velocity $v_{r}$. The amount of ICM mass swept per unit time by the disk is
\begin{equation*}
    \dot{M} = \pi R_{d}^{2}\rho_{\text{ICM}}v_{r}
\end{equation*}
If we assume that the wind is stopped by the interstellar medium (ISM) of the disk, then the momentum transferred to the disk per unit time (that is a force) is
\begin{equation*}
    \dot{p} = \dot{M}v_{r} =  \pi R_{d}^{2}\rho_{\text{ICM}}v_{r}^{2}
\end{equation*}
The pressure is just 
\begin{equation*}
    P_{\text{ram}} = \rho_{\text{ICM}}v_{r}^{2}
\end{equation*}
The gravitational field (acceleration) due to the disk is approximately $2\pi G \Sigma_{*}$ and so the gravitational force per unit area (pressure) on the interstellar gas is $2\pi G \Sigma_{*}\Sigma_{\text{gas}}$.
\par Finally, we expect that stripping occurs if
\begin{equation}
    P_{\text{ram}} = \rho_{\text{ICM}}v_{r}^{2} > 2\pi G \Sigma_{*}\Sigma_{\text{gas}}
    \label{eq:rps}
\end{equation}
\par According to tuned models and hydrodynamic simulations of \emph{ram pressure stripping} (RPS), the process removes the various components of a galaxy’s ISM starting from its outer regions and progressing inward.
\par Whereas in dwarf systems the stripping is generally complete, in massive systems some gas can be retained in the inner region, where the gravitational potential well of the galaxy is the deepest. The activity of star formation is reduced by the lack of gas (strangulation), and the perturbed galaxies thus migrate below the main sequence becoming quiescent, red objects without affecting their kinematic properties and the structure of old stellar populations.
\par This quenching process is rapid ($<1\unit{Gyr}$), in particular in dwarf system where most of the gas is removed at the first infall of the galaxy within the ICM before the pericenter of their orbit.
\par Only under particular configurations such as edge-on stripping, when the perturbed gas is pushed along the disk, there might be a short lived burst of the star formation activity of the perturbed galaxies. Star formation can be present in the tail of stripped material, but this does not occur in all RPS systems.
\begin{figure}[h!]
    \centering 
    \includegraphics[width=0.7\textwidth]{img/jellyfishgal.png}
    \caption{D100, a barred spiral galaxy close to the centre of the Coma cluster. Jellyfish galaxies represent an extreme example of ram-pressure stripping, where the stripped gas streams out in a long tail behind the galaxy, giving them their distinctive look. The HST observations found 37 bright patches, and analysis of their colors showed 10 of them to be clumps of star formation, all of which are found in the tail of gas. The 27 other sources are mostly background sources, such as distant galaxies. Credits: \cite{Cramer_2019}.}
\end{figure}

\newpage
\section{Collisionless Dynamics}

If star-star collisions are negligible, stars move in the smooth gravitational field $\Phi(x,t)$ of the galaxy. In this case a galaxy may be described by a probability density in phase-space called \emph{distribution function}. We define a distribution function $f = f(\vec{x}, \vec{v}, t)$ so that the number of particles in range $[\vec{x}, \vec{x}+\text{d}\vec{x}\,]$ and $[\vec{v}, \vec{v}+\text{d}\vec{v}\,]$ is given by 
\begin{equation}
    \text{d}N(t) = f(\vec{x}, \vec{v}, t)\,{d}\vec{x}\,\text{d}\vec{v}
\end{equation}
To be properly defined, $f$ must meet the usual requirements for distribution functions
\begin{equation}
    f \geqslant 0 \quad \int f\,\text{d}\vec{x} < +\infty \quad \int f\,\text{d}\vec{v} < +\infty
\end{equation}
Thanks to these properties we can easily define the number and mass density 
\begin{equation*}
    n(\vec{x}, t) = \int f(\vec{x}, \vec{v}, t)\,\text{d}\vec{v}
\end{equation*}
as well as the concept of \emph{averaged quantities}, for example the average velocity
\begin{equation*}
    \langle \vec{v}\,\rangle = \frac{1}{n}\int \vec{u}f(\vec{x}, \vec{u}, t)\,\text{d}\vec{u}
\end{equation*}
\par Liouville's theorem\footnote{The proper formulation would be that of a continuity equation for the states in the phase space $\partial_{t}\rho +\nabla(\rho\vec{v}) = 0$.} tells us that the six dimensional volume in phase space is conserved $\text{d}^{3}x\text{d}^{3}v \approx \text{d}^{3}x_{0}\text{d}^{3}v_{0}$ up to first order\footnote{Properly speaking, we'd have to consider the transformation $\vec{x} = \vec{x}_{0}+\vec{u}_{0}\,\text{d}t$, $\vec{v} = \vec{v}_{0}+\vec{a}_{0}\,\text{d}t$, which has a Jacobian $\left\|J\right\| = 1+\text{o}(\text{d}t^{2})$.} in $\text{d}t$. This way we can write the initial infinitesimal number of particles in volume $\text{d}^{3}x_{0}\text{d}^{3}v_{0}$ and the number of particles in volume $\text{d}^{3}x\text{d}^{3}v$
\begin{align*}
    \text{d}N_{0} &= f(\vec{x}_{0}, \vec{v}_{0}, t)\,\text{d}^{3}x_{0}\text{d}^{3}v_{0}\\
    \text{d}N &= f(\vec{x}, \vec{v}, t)\,\text{d}^{3}x\text{d}^{3}v
\end{align*}
Requiring the number of particles to be conserved $\text{d}N_{0} = \text{d}N$, we end up with 
\begin{equation*}
    f(\vec{x}_{0}, \vec{v}_{0}, t) = f(\vec{x}_{0}+\vec{u}_{0}\,\text{d}t, \vec{v}_{0}+\vec{a}_{0}\,\text{d}t, t) \approx f(\vec{x}_{0}, \vec{v}_{0}, t) +\frac{\text{d}f}{\text{d}t}
\end{equation*}
which implies $\text{d}f/\text{d}t = 0$. Writing down what this implies component by component we get 
\begin{equation}
    \boxed{
    \partial_{t}f + v^{i}\partial_{x^{i}}f+a^{i}\partial_{v^{i}}f = 0
    }
    \label{eq:vlasov}
\end{equation}
Note that sometimes it might be useful to express the $i$-th component of the acceleration as the gradient of the external potential $a^{i} = \partial^{i}\Phi$.
This is called \emph{Collision-less Boltzmann equation} (CBE). Please, don't bother to ask why of the weird indices placement. Whichever way you put the indices is (here) absolutely irrelevant.
\par The flow of stellar phase points through phase space is incompressible; if you follow a star along its orbit, the density in the phase space (6D) around the star is constant.
\par If there are close encounters between stars, these can alter the position and velocity much faster than a smoothed potential. In this case, the effects are given as an extra “collisional” term on the right side.
\par So far we've been ignoring the possibility of Liouville's theorem being violated, that is new states could be popping out or some could be disappearing. To address this possibility, we have to add a formation rate (i.e., the SFR) and a "death rate" for stars, thus modifying the CBE as such
\begin{equation*}
    \partial_{t}f + v^{i}\partial_{x^{i}}f+a^{i}\partial_{v^{i}}f = \text{SFR}-\text{D}
\end{equation*}
The CBE would be vaild only if
\begin{equation*}
    \text{SFR}-\text{D} \ll \frac{f}{\tau_{\text{cross}}}
\end{equation*}
which is a condition that is generally satisfied in all elliptical galaxies: $\text{SFR}\approx 0$ and $\text{D}\approx 0$, for there are only slow evolving galaxies. For star forming galaxies, on the other hand, the CBE starts to be not so accurate.

\subsection{Jeans Equations}

The collisionless Boltzmann equation is rather useless\footnote{Poor Boltzmann.}, as it is not only a function of 7 variables (many of which are not easily observable), but of their derivatives as well. Even if $f$ were measurable, there would be large uncertainties due to Poisson noise. Fortunately, more tractable equations can be found by finding the momenta of the equation. Such momenta produce the Jeans Equations.
\par I'm going to skip explicit calculations because they're somewhat long, tedius, and there's not much really to learn from making that much algebra. The only things you have to remember are the definitions of the star density $n$ as the integral of the distribution function over velocities, the definition of averaged quantity (still over velocities) and the fact that position and speed are independent (since they're conjugate variables).
\par Integration with respect to velocities leads to the first Jeans equation
\begin{equation}
    \dot{n}+\partial^{i}(n\langle v_{i}\rangle) = 0
    \label{eq:jeans_1st}
\end{equation}
Then again, we can multiply the CBE by a factor $v_{j}$ and then integrate. At the end of it, you'll find the second Jeans equation (which is actually a vectorial one)
\begin{equation}
    \partial_{t}(n\langle v_{j}\rangle ) + \partial^{i}\left(n\langle v_{j}v_{i}\rangle \right)+\partial^{j}\Phi n = 0
    \label{eq:jeans_2nd}
\end{equation}
The final Jeans equation comes from observing that the anisotropic pressure term (i.e., the stress tensor), $\sigma_{ij}$ describes the random stellar motions (velocity dispersion or stress tensor)
\begin{equation}
    \sigma_{ij}^{2} = \langle v_{j}v_{i}\rangle - \langle v_{j}\rangle\langle v_{i}\rangle
    \label{eq:anisotropic_pressure}
\end{equation}
which we can subsitute into the second Jeans equation
\begin{equation}
    \partial_{t}(n\langle v_{j}\rangle ) + \partial^{i}\left(n( \sigma_{ij}^{2} + \langle v_{j}\rangle\langle v_{i}\rangle ) \right)+n\,\partial^{j}\Phi = 0
    \label{eq:jeans_anisotropic}
\end{equation}
Here, $n\sigma_{ij}^{2}$ is the anisotropic pressure (stress) tensor. But since $\sigma_{ij}^{2}$ is symmetric, its matrix can be diagonalized. This produces the principal axes of the velocity ellipsoid. However, pressures in stellar systems can still be anisotropic $\sigma_{11}\neq \sigma_{22}\neq \sigma_{33}$.
\par The incompressible stellar fluid experiences anisotropic pressure forces. If $\Phi$ and $n$ are known, the Jeans equations have 9 unknowns (3 streaming motions and at least 3 terms of the random motion tensor). But we have only four equations, meaning this is not a closed set.
\par For comparison, in fluid dynamics there are only 4 unknowns (3 streaming motions and the pressure). The 3 Euler Equations combined with the Equation of State forms a closed set. Fluids are isotropic . This difference is due to the collisional nature of fluids, that are characterized by an isotropic velocity dispersion tensor (i.e. the velocity ellipsoid becomes a sphere).

\subsubsection*{Applications}

Jeans equations can be written for cylindrically symmetric systems. Consider the equation describing the $z$ direction
\begin{equation}
    \partial_{t}(n\overline{v}_{z})+\partial_{r}(n\overline{v_{r}v_{z}})+\partial_{z}(n\overline{v_{z}^{2}})+\frac{n\overline{v_{r}v_{z}}}{r} + n\partial_{z}\Phi = 0
    \label{eq:jeans_cylindrical}
\end{equation}
and assume now a steady state system. The former equation turns into
\begin{equation*}
    \frac{1}{r}\partial_{r}(Rn\overline{v_{r}v_{z}})+\partial_{z}(n\overline{v_{z}^{2}}) = - n\partial_{z}\Phi
\end{equation*}
We are interested in the density of a thin disk, where the density falls off much faster in $z$ than in $r$. Typically, disks are a few $100\unit{pc}$ thick, with a radial scale of a few kpc, so we can neglect the radial derivative. We're left with 
\begin{equation*}
    \partial_{z}(n\overline{v_{z}^{2}}) = - n\partial_{z}\Phi
\end{equation*}
which simply states that vertical pressure is balancing vertical gravity (hydrostatic equilibrium). This is the Jeans equation for a one-dimensional slab. We can also show that Poisson's equation in a thin disk approximation is 
\begin{equation*}
    \parparder{\Phi}{z} = 4\pi G \rho
\end{equation*}
with $\rho$ the \emph{total} density (stars, gas and DM).
\par We can differentiate one more time our Jeans equation and find 
\begin{equation*}
    \partial_{z}\left(\frac{1}{n}\partial_{z}(n\overline{v_{z}^{2}}) \right) = -4\pi G \rho
\end{equation*}
The $n$ here is the number density of G stars or whatever type is chosen. Thus if for any population of stars we can measure $\overline{v_{z}^{2}}$ and $n$ as a function of height $z$, we can calculate the total local density $\rho$. Using this technique for F and K giants stars, Oort found
\begin{equation*}
    \rho_{0} = \rho(R_{0}, z=0) = 0.15M_{\odot}\unit{pc}^{-3}
\end{equation*}
while more recent estimates set it to
\begin{equation}
    \rho_{0} = \rho(R_{0}, z=0) = 0.10 \pm 0.01M_{\odot}\unit{pc}^{-3}
    \label{eq:oortlimit}
\end{equation}
Please note that this is \emph{not} the visible mass, but the \emph{total} inferred mass from the procedure we've (well, Oort has) devised. Over a scale height of a few hundreds pc, there is little evidence of DM.
\par Note that we can determine instead more accurately (since there is one less difference, or differential, involved) the surface density
\begin{equation*}
    \Sigma(z) = \int_{-z}^{z}\rho(z)\diff{z} = -\frac{1}{2\pi G n}\partial_{z}(n\overline{v_{z}^{2}})
\end{equation*}
which a dynamical estimate sets at
\begin{equation*}
    \Sigma(R_{0}, 1.1\unit{kpc}) \approx 71 \pm 6 M_{\odot}\unit{pc}^{-2}
\end{equation*}
for which Baryons account of only about half (includes stars and gas)
\begin{equation*}
      \Sigma(R_{0}, 1.1\unit{kpc})_{B} \approx 41 \pm 15 M_{\odot}\unit{pc}^{-2}
\end{equation*}
Hence at $R_{0}$ the dark halo contribution to the vertical surface density becomes significant only when the column is extended to large enough $z$, because the halo is diffuse and vertically extended.

\subsubsection*{Spherically Symmetric Steady State Systems}

The expression is
\begin{equation}
    \frac{1}{n}\partial_{r}(n\sigma_{rr}^{2})+\frac{2(\sigma_{rr}^{2}-\sigma_{t}^{2})}{r} = -\partial_{r}\Phi
    \label{eq:jeans_spherical}
\end{equation}
It is useful to define the anisotropy parameter
\begin{equation}
    \beta(r) = 1-\frac{\sigma_{t}^{2}}{\sigma_{rr}^{2}}
\end{equation}
For the spherical simmetry we also have
\begin{equation*}
    -\partial_{r}\Phi = -\frac{GM(r)}{r^{2}}
\end{equation*}
so that 
\begin{equation}
    M(<r) = -\frac{r\sigma_{rr}^{2}}{G}\left[\der{\ln n}{\ln r}+\der{\ln\sigma_{rr}^{2}}{\ln r}+2\beta(r)\right]
\end{equation}
Thus, if we can measure $n(r)$, $\beta(r)$, $\sigma_{rr}$, we can use the spherical Jeans equation to infer the mass profile $M(r)$.
\par Consider an external, spherical galaxy. Observationally, we can measure the projected surface brightness profile $I(r)$ (from which, if we assume a sensible $M/L$, we can derive a projected two-dimensional stellar density $\Sigma(r)$).
\par If $I(r)$ is circularly symmetric, it is possible that the luminosity density $J(r)$ is spherically symmetric as well. So
\begin{equation*}
    I(R) = \int_{-\infty}^{\infty}J(r)\diff{z} = 2\int_{R}^{\infty}\frac{J(r)}{\sqrt{r^{2}-R^{2}}}r\diff{r}
\end{equation*}
which can be inverted since it's an Abel integral equation
\begin{equation*}
    J(r) = -\frac{1}{\pi}\frac{\text{d}}{\text{d}r}\int_{r}^{\infty}I(R)\frac{1}{\sqrt{R^{2}-r^{2}}}\diff{R}
\end{equation*}
Then, using a $M/L$ we can transform $J(r)$ to $\rho(r)$ (and $n(r)$ assuming that all stars have the same mass or a specific IMF).
\par Similarly, the line-of-sight velocity dispersion is an observationally accessible quantity
\begin{equation*}
    \Sigma(R)\sigma_{\text{LOS}}^{2}(R) = 2\int_{R}^{\infty}\left(1-\beta\frac{R^{2}}{r^{2}}\right)\frac{\langle v_{r}^{2}\rangle n(r)r}{\sqrt{r^{2}-R^{2}}}\diff{r}
\end{equation*}
Unfortunately we have only one observable $\sigma_{\text{LOS}}$, so we have to do guesses on either $\langle v_{r}^{2}\rangle$ or $\beta$. For example, if $\beta = 0$ we could use the Abel inversion to find $\langle v_{r}^{2}\rangle$.
\par This formalism has a useful application to galaxy clusters, since 
\begin{equation*}
    M(<r) = kr\sigma_{rr}^{2}(r)
\end{equation*}
where $k = 3\pm 1.5$ depending on the particular details. We can do a order of magnitude estimate
\begin{equation*}
    \frac{M_{\text{cluster}}}{N_{\text{gal}}M_{\text{gal}}} \approx 10
\end{equation*}
thus the need to assume the presence of DM.

\section{Elliptical Galaxies}

If the subject is particularly of your liking, consider giving a read to \cite{Kormendy_2009}, for they have written a rather extensive review about the structure and formation of elliptical and spheroidal galaxies\footnote{Yeah, totally random choice of words on my side.}.

\subsection{Generalities}

Unlike disk galaxies, which we've briefly discussed in earlier chapters, elliptical galaxies have surface brightness profiles that are well described by a Sérsic profile
\begin{equation}
    I(R) = I_{e}\exp\left[-b_{n}\left(\left(\frac{R}{R_{e}}\right)^{1/n}-1\right)\right]
    \label{eq:sersicprofile}
\end{equation}
where within effective radius $R_{e}$ is contained half the light of the galaxy. For $n=4$ it is equivalent to the de Vaucouleurs profile (\ref{eq:devaucouleurs}) we've already encountered.
\par Elliptical galaxies of different luminosity have different concentrations (i.e., different $n$-indices). Using the Sérsic profile (\ref{eq:sersicprofile}) to fit for the $n$-index gives $n<4$ for low luminosity elliptical (which are less concentrated), while $n>4$ for higher luminosity ellipticals.
\par Elliptical galaxies are broadly described by Sérsic profiles, but their innermost regions show a dichotomy. Massive ellipticals often have central light deficits, or \emph{cores}, with respect to the outer Sérsic profile. Lower-luminosity ellipticals are usually coreless or \emph{cuspy}, with surface brightness profiles that continue rising steeply toward the center. In particular, the more luminous ellipticals ($M_{V}\lesssim -21.7$) tend to present cores, inducing a flattening in the central region of the surface brightness profile. On the other hand, midsize ellipticals ($-21.5\lesssim M_{V}\lesssim -15.5$) are typically core-less systems where the SB profile steeply rises to center.

\subsection{Faber-Jackson Relation}

Just as for spiral galaxies, stars move faster in more luminous galaxies, following a trend roughly predicted by the Tully-Fisher relation (\ref{eq:tullyfisher}). At the centers of bright ellipticals, the dispersion can reach $500 \unit{km}\unit{s}^{-1}$ while $\sigma \sim 50\unit{km}\unit{s}^{-1}$ in the least luminous objects. In the V band, it has been found the following empirical relation
\begin{equation}
    \frac{L_{V}}{2\power{10}{10}L_{\odot}} \approx \left(\frac{\sigma}{200 \unit{km}\unit{s}^{-1}}\right)^{4}
    \label{eq:faberjackson}
\end{equation}
which is often called the Faber-Jackson relation, and is essentially the equivalent of the Tully-Fisher relation for elliptical galaxies.

\subsection{Tensor Virial Theorem}

As we'll see in the following, when dealing with elliptical galaxies, a crucial distinction has to be made. Are ellipticals deformated because they're just strongly asymmetric or is it because they're strongly rotating?
\par To find an answer to this question, we're going to need the \emph{tensor virial theorem}. For this, I'll just cover the definitions and then skip right to the final result (\emph{Said Lying Odysseus}).
\par Consider the second Jeans equation (\ref{eq:jeans_2nd}) and multiply it by $x_{k}$, then integrate over the spatial coordinates. Doing so, we obtain 
\begin{equation}
    \int x_{k}\partial_{t}(n\langle v_{j}\rangle)\,\text{d}^{3}x = -\int x_{k}\partial^{i}(n\langle v_{i}v_{j}\rangle)\,\text{d}^{3}x - \int x_{k}\partial_{j}\Phi n\,\text{d}^{3}x
    \label{eq:starting_tvt}
\end{equation}
We define Chandrasekhar's potential tensor from the last term of the former equation 
\begin{equation}
    W_{kj} = -\int x_{k}n\partial_{j}\Phi\,\text{d}^{3}x
    \label{eq:chandrapot}
\end{equation}
Since by definition the gravitational potential (per unit mass) is given by 
\begin{equation*}
    \Phi(\vec{x}) = -G\int \frac{n(\vec{x}')}{|\vec{x}-\vec{x}'|}\,\text{d}^{3}x'
\end{equation*}
we can plug it into (\ref{eq:chandrapot}) and find, with some algebra that exploits the symmetry $x\longleftrightarrow x'$, the following expression 
\begin{equation*}
    W_{kj} = -\frac{G}{2}\left[\int\int n(\vec{x}')n(\vec{x})\frac{(x_{j}'-x_{j})(x_{k}'-x_{k})}{|\vec{x}'-\vec{x}|^{3}}\right]\,\text{d}^{3}x\,\text{d}^{3}x'
\end{equation*} 
which is also manifestly symmetric. Taking the trace $W^{k}_{\;\;k}$, we find out that it reduces to the scalar gravitational potential $W$
\begin{equation*}
    W^{k}_{\;\;k} = -\frac{G}{2}\left[\int\int n(\vec{x}')n(\vec{x})\frac{1}{|\vec{x}'-\vec{x}|}\right]\,\text{d}^{3}x\,\text{d}^{3}x' = \frac{1}{2}\int n\Phi\,\text{d}^{3}x = W
\end{equation*}
where we've simply subbed in the definition of $\Phi$.
\par Moving to the second term of (\ref{eq:starting_tvt}), we can integrate per parts, apply the divergence theorem and find 
\begin{equation*}
    -\int x_{k}\partial^{i}(n\langle v_{i}v_{j}\rangle)\,\text{d}^{3}x = \int n\langle v_{i}v_{j}\rangle\,\text{d}^{3}x
\end{equation*}
We define the \emph{symmetric kinetic tensor}
\begin{equation}
    K_{jk} = \frac{1}{2}\int n \langle v_{i}v_{j}\rangle\,\text{d}^{3}x
    \label{eq:skt}
\end{equation}
which we split into the contribution of ordered ($T_{jk}$) and random ($\Pi_{jk}$) motions in the following way\footnote{Recall that $\sigma_{jk}^{2} = \langle v_{j}v_{k}\rangle - \langle v_{j}\rangle \langle v_{k}\rangle$.}
\begin{equation}
    T_{jk} = \frac{1}{2}\int n \langle v_{j}\rangle\langle v_{k}\rangle\,\text{d}^{3}x \quad \Pi_{jk} = \int n \sigma^{2}_{jk}\,\text{d}^{3}x
    \label{eq:kinetictensor_decomp}
\end{equation}
so that 
\begin{equation*}
    K_{jk} = T_{jk}+\frac{1}{2}\Pi_{jk}
\end{equation*}
Equation (\ref{eq:starting_tvt}) then becomes 
\begin{equation*}
    \int x_{k}\partial_{t}(n\langle v_{j}\rangle)\,\text{d}^{3}x = W_{jk}+2T_{jk}+\Pi_{jk}
\end{equation*}
This is symmetric (because everything on the right is symmetric), so exchanging the ($k$, $j$) and the ($j$, $k$) components and averaging
\begin{equation*}
    \frac{1}{2}\left[\int x_{j}\partial_{t}(n\langle v_{k}\rangle)\,\text{d}^{3}x + \int x_{k}\partial_{t}(n\langle v_{j}\rangle)\,\text{d}^{3}x \right] = W_{jk}+2T_{jk}+\Pi_{jk}
\end{equation*}
The partial time derivatives on the left-hand side can be pulled outside the integral and then changed into a complete derivative (because the spatial part is integrated out)
\begin{equation*}
    \frac{1}{2}\frac{\text{d}}{\text{d}t}\left[\int x_{j}(n\langle v_{k}\rangle)\,\text{d}^{3}x + \int x_{k}(n\langle v_{j}\rangle)\,\text{d}^{3}x \right] = W_{jk}+2T_{jk}+\Pi_{jk}
\end{equation*}
To define this expression, it is useful to define the symmetric moment of inertia tensor and its temporal derivative 
\begin{equation}
    I_{jk} = \int n x_{j}x_{k}\,\text{d}^{3}x \quad \der{I_{jk}}{t} = \int \parder{n}{t}x_{j}x_{k}\,\text{d}^{3}x
    \label{eq:inertiatensor}
\end{equation}
Using the first Jeans equation (\ref{eq:jeans_1st}) and the divergence theorem, we obtain 
\begin{equation*}
    \der{I_{jk}}{t} = \int n(x_{j}\langle v_{k}\rangle+ x_{k}\langle v_{j}\rangle )\,\text{d}^{3}x
\end{equation*}
so that we can finally present the Tensor Virial Theorem
\begin{equation}
    \boxed{
    \frac{1}{2}\dder{I_{jk}}{t} = W_{jk}+2T_{jk}+\Pi_{jk}
    }
    \label{eq:tensorvirialtheorem}
\end{equation}
If the system is in a steady-state, then the moment of inertia tensor is stationary, and the Tensor Virial Theorem reduces to
\begin{equation*}
    0 = W_{jk}+2K_{jk}
\end{equation*}
whose trace reduces to the familiar scalar virial theorem (\ref{eq:virial_theorem}).

\subsection{Flattening of Elliptical Galaxies}

Consider now the case of an axisymmetric galaxy that rotates around its symmetry axis (say, the $z$ axis) and is seen edge-on. The $x$ axis is taken to be parallel to the line of sight, so that the $yz$ plane is actually the sky plane. Symmetry sentences the death of all non-diagonal terms of $W_{jk}$, $T_{jk}$, $\Pi_{jk}$. Rotational invariance on the plane also requires $xx$ terms to be equal to $yy$ terms. We're then left with just two equations (the third is not independent)
\begin{align*}
    2T_{xx}+\Pi_{xx} &= -W_{xx}\\
    2T_{zz}+\Pi_{zz} &= -W_{zz}
\end{align*}
of which we can take the ratio
\begin{equation*}
    \frac{2T_{xx}+\Pi_{xx}}{2T_{zz}+\Pi_{zz}} = \frac{W_{xx}}{W_{zz}}
\end{equation*}
Note that the right hand side depends only on the geometry of the system. Let us consider for example an isotropic, oblate rotator. Isotropy requires $\Pi_{xx} = \Pi_{zz} = M\sigma^{2}$ and $T_{zz} = 0$, while the trace of the $T$ tensor is 
\begin{equation*}
    T^{i}_{\;\;i} = 2T_{xx} = \frac{1}{2}\int n(\vec{x})(\langle v_{\phi}\rangle^{2}\sin^{2}\phi + \langle v_{\phi}\rangle^{2}\cos^{2}\phi)\,\text{d}^{3}x := \frac{1}{2}Mv^{2}_{\text{rot}}
\end{equation*}
where $v_{\text{rot}}$ is the density weighted rotation speed.
\par The ratio we've previously computed then becomes 
\begin{equation*}
    \frac{\frac{1}{2}Mv^{2}_{\text{rot}}+M\sigma^{2}}{M\sigma^{2}} = \frac{W_{xx}}{W_{zz}}
\end{equation*}
It can be shown that the right hand side of the equation, the geometrical part, can be expressed as the ratio of the major axis of the ellipsoid or as a function of the ellipticity $\epsilon = 1-b/a$, finding 
\begin{equation}
    \frac{v_{\text{rot}}}{\sigma} \approx \sqrt{2\left[(1-\epsilon)^{-0.9}-1\right]} = \sqrt{2\left[\left(\frac{b}{a}\right)^{-0.9}-1\right]}
    \label{eq:isoblator}
\end{equation}
This specifies the relation between the flattening of the spheroid and the ratio of streaming motion to random motion. Also, note that you need a rather large amount of rotation to achieve only modest flattening. In fact, according to this relation, even fairly round galaxies should rotate quite fast; for example
\begin{equation*}
    \frac{b}{a} = 0.7 \implies \frac{v_{\text{rot}}}{\sigma} \approx 0.9
\end{equation*}
which is an absurd value.
\par Consider then a non-rotating, anisotropic, oblate system\footnote{The system is still axisymmetric, the anisotropy is limited to the velocity fluctuations along the $z$ axis.}. The velocity dispersion along the $x$ and $z$ axis will then be different, and repeating the calculations we'd find 
\begin{equation}
    \frac{\sigma_{zz}}{\sigma_{xx}} \approx \left(\frac{b}{a}\right)^{0.45}
    \label{eq:anioblnonrot}
\end{equation}
Now, the same 0.7 flattening requires only a small anisotropy $\sigma_{zz}/\sigma_{xx}\approx 0.85$.
\par Finally, let us consider a rotating, anisotropic, oblate system, for which the anisotropy is characterized by a $\delta$ parameter. We can write 
\begin{align*}
    &\Pi_{zz} = (1-\delta)\Pi_{xx} = (1-\delta)M\sigma^{2}\\
    & T_{zz} = 0 \quad 2T_{xx} = \frac{1}{2}Mv^{2}_{\text{rot}}
\end{align*}
from which 
\begin{equation}
    \frac{v_{\text{rot}}}{\sigma} \approx \sqrt{2\left[(1-\delta)(1-\epsilon)^{-0.9}-1\right]}
    \label{eq:anisoblator}
\end{equation}
\par A potential problem is that we can not directly measure $v_{\text{rot}}$ nor $\sigma$. Rather, we measure properties that are projected along the line-of-sight. Furthermore, in general we don’t see a system edge-on but under some unknown inclination angle $i$ (so, we do not see the full major axis, but just their projection).
\begin{figure}[h!]
    \centering 
    \includegraphics[width=0.7\textwidth]{img/elligalaxiesrot.png}
\end{figure}
\par Elliptical galaxies can thus be split into two kinematic classes
\begin{itemize}
    \item \emph{Fast rotators}: Kinematics consistent with oblate rotators, their shape is flattened by rotation;
    \item \emph{Slow rotators}: They present very little rotation, and their shape is due to anisotropic pressure support.
\end{itemize}
\par The isophotes\footnote{"In geometry, an isophote is a curve on an illuminated surface that connects points of equal brightness."} of elliptical galaxies are elliptical, but not perfectly so. They show deviations that are typically either '\emph{disky}' or '\emph{boxy}', as shown in Fig.\ref{fgr:diskyboxy}.
\begin{figure}[h!]
    \centering
    \includegraphics[width=0.7\textwidth]{img/diskyboxy.png}
    \caption{On the left, a boxy galaxy; on the right, a disky galaxy.}
    \label{fgr:diskyboxy}
\end{figure}
\par The angular variation in the radial coordinate can be expressed in terms of a Fourier series
\begin{equation}
    \Delta R(\phi) = R_{\text{isop}}(\phi) - R_{\text{best-ell}}(\phi) = \sum_{n=1}^{\infty}\left[a_{n}\cos(n\phi)+b_{n}\sin(n\phi)\right]
\end{equation}
Disky and boxy galaxies correspond, respectively, to FT expansions with $a_{4}>0$ and $a_{4}<0$. As much as 80\% of elliptical galaxies shows deviations from purely elliptical isophotes.
\begin{figure}[h!]
    \centering
    \includegraphics[width=0.45\textwidth]{img/disky.png}
    \includegraphics[width=0.45\textwidth]{img/boxy.png}
    \caption{On the left, R-image of NGC4660, an elliptical galaxy with a disk component of the isophotes; on the right, R-image of NGC5322, an elliptical galaxy with box-shaped isophotes.}
    \label{fgr:disky_boxy}
\end{figure}

\subsection{Isophotal Twist}

Some ellipticals galaxies reveal isophotal twists, with direction of major axis of isophote changing with isophotal level.
\par The simplest explanation is that elliptical galaxies can be triaxial (rather than oblate or prolate), and have their intrinsic axis ratios change with radius. Such a system in projection will reveal an \emph{isophotal twist}.
\par As a rule of thumb, you might want to remember that luminous ellipticals are usually more boxy and twisted.
\par The disky ellipticals may form from mergers of gas-rich disk galaxies; the merger remnants will then contain a population of young stars that were formed in the process of merging, and depending on the relative orientation of the angular momentum vectors of the two galaxies, the merged galaxy may have a substantial contribution of rotational support.
\par On top of that, core-less galaxies also reveal some "extra light" which might be the result of nuclear starburst resulting from "wet mergers" (with gas), while galaxies with cores could be the result of dry mergers that created a more diffuse central nucleus.

\subsection{Integrals of Motion: Jeans Theorem}

Let’s consider a system in steady state, where $\Phi$ and the distribution function $f$ are not explicit functions of time. In this
case, we may introduce the integrals of motion. An "integral of motion" is a function $I(\vec{r}, \vec{v})$ (not explicit function of time) which is constant along a star's orbit. Obvious examples of possible integrals of motion are the energy (in a static potential), the whole angular momentum (in a static potential with spherical symmetry), the off-plane component of the angular momentum (in an axisymmetric, static potential).
\par Since $I$ is constant along an orbit, we can expand its time derivative
\begin{equation*}
    \der{I}{t} = \partial^{i}I\der{x_{i}}{t} + \parder{I}{v_{i}}\der{v_{i}}{t} = (\vec{v}\cdot \nabla)I-\nabla\Phi\parder{I}{\vec{v}} = 0 = \text{CBE}(I)
\end{equation*}
where the last two equalities hold if $I$ really is an integral of motion. What this implies is truly remarkable: $I$ is a solution of the collisionless Boltzmann equation (\ref{eq:vlasov}) and it's actually a valid distribution function.
\par Jeans' theorem proves that any steady-state solution of the collisionless Boltzmann equation depends on the phase-space coordinates only through integrals of motion in the given potential, and any function of the integrals yields a steady-state solution of the collisionless Boltzmann equation. In short
\begin{equation}
    \der{f}{t}[I_{1}(\vec{r},\vec{v}),...,I_{n}(\vec{r},\vec{v})] = \parder{f}{I_{i}}\der{I_{i}}{t}
    \label{eq:jeanstheorem}
\end{equation}
\par The Jeans theorem is extremely useful since it allows us to construct legitimate distribution functions $f$ using integrals of motion. In the special case of steady state spherical systems, it can be shown that:
\begin{enumerate}
    \item Distribution functions must be of the form $f(E, |L|)$;
    \item Distribution functions of the form $f(E)$ must have an isotropic velocity dispersion $\sigma_{r} = \sigma_{\theta} = \sigma_{\phi}$;
    \item Distribution functions of the form $f(E, |L|)$ must have an anisotropic velocity dispersion $\sigma_{r} \neq \sigma_{\theta}=\sigma_{\phi}$.
\end{enumerate}
\par Clearly, an important step is to require that the distribution function also yields the potential $\Phi(r)$
\begin{equation*}
    4\pi G \int f(\vec{r},\vec{v})\,\text{d}^{3}v = 4\pi G n = \nabla^{2}\Phi
\end{equation*}
which in spherical coordinates becomes 
\begin{equation*}
    \frac{1}{r^{2}}\frac{\text{d}}{\text{d}r}\left(r^{2}\der{\Phi}{r}\right) = 4\pi G \int f\left(\frac{1}{2}v^{2}+\Phi, |\vec{r}\wedge\vec{v}|\right)\,\text{d}^{3}v
\end{equation*}
Solutions of this equation give self-consistent spherical stellar systems, provided the distribution function is non-negative and the resulting density and potential are physically acceptable.
\par To make things easier we redefine the potential and the energy by adjusting the arbitrary constant and changing the sign
\begin{align*}
    \Psi &= -\Phi +\Phi_{0}\\
    \epsilon &= -E +\Phi_{0} = \Psi -\frac{1}{2}v^{2}
\end{align*}
We choose $\Phi_{0}$ such that $f>0$ for $\epsilon>0$ (bound orbits) and $f=0$ for $\epsilon \leq 0$ (unbound orbits). The iter is for example the following:
\begin{enumerate}
    \item  Choose a distribution function which is a function of energy $f(\epsilon)$. Because of Jeans Theorem, $f(E)$ is already a valid solution to the steady-state CBE;
    \item Integrate the distribution function over $v$ to find $n(\Psi)$;
    \item Solve Poisson’s equation to find $\Psi(r)$;
    \item Plug $\Psi(r)$ into $n(\Psi)$ to find the mass distribution $n(r)$.
\end{enumerate}
Consider for example an isolated, self-gravitating, steady-state stellar system described by
\begin{equation*}
    \frac{1}{r^{2}}\frac{\text{d}}{\text{d}r}\left(r^{2}\der{\Psi}{r}\right) = -16\pi^{2} G \int f\left(\epsilon\right)\sqrt{2(\Psi-\epsilon)}\,\text{d}\epsilon
\end{equation*}
and take a polytropic sphere with a power law distribution function $f(\epsilon) = F\epsilon^{n-3/2}$ for $\epsilon > 0$.
\par Following our charted iter
\begin{align*}
    n &= 4\pi \int_{0}^{\infty} f(\epsilon)v^{2}\diff{v}\\
    &= 4\pi F\int_{0}^{\infty}\left(\Psi - \frac{1}{2}v^{2}\right)^{n-3/2}v^{2}\diff{v}
\end{align*}
which is solved by $n(\Psi) = c_{n}\Psi^{n}$. Subbing it in the spherical Poisson equation retrieves the Lane-Emden equation. Thus, we find that for a self-gravitating sphere, the density profile $\rho(r)$ is the same for stars with distribution function $\epsilon^{n-3/2}$ and gases with polytropic EoS.
\par For $n<5$, the Lane-Emden equation admits solutions with vanishing density at a finite radius\footnote{Systems with $n>5$ have infinite mass at large $r$, while for systems with $n=\infty$ (isothermal equation) the analysis we've given breaks down.}, but the analitical solutions are limited to only a few values of $n$, such as $n=5$ (Plummer Sphere).
\par For $n=5$ the density profile is 
\begin{equation*}
    \rho \sim \left(1+\left(\frac{r}{b}\right)^{2}\right)^{-5/2}
\end{equation*}
which gives a finite mass, is well-behaved at $r=0$ and is a good match for the density profile of globular clusters and dwarf galaxies.
\par Consider an isothermal sphere described by
\begin{equation*}
    f(\epsilon)\sim \frac{\rho_{1}}{(2\pi \sigma^{2})^{3/2}}e^{\epsilon/\sigma^{2}}
\end{equation*}
\par More positive $\epsilon$ means more bound. Also, note that $f(\epsilon) > 0$ for $\epsilon <0$, so the equation admits unbound stars. Following the usual procedure gives $\rho = \rho_{1}\exp(\Psi/\sigma^{2})$, and plugging into Poisson's equation we find a solution in the \emph{singular isothermal sphere} (SIS)
\begin{equation*}
    n = \frac{\sigma^{2}}{2\pi G r^{2}}
\end{equation*}
which has two rather nice properties (circular velocities and velocity dispersions independent of radius) and two not-so-nice properties (infinite density at $r=0$ and infinite mass for large radii).
\par The general isothermal sphere takes as boundary conditions at $r=0$
\begin{equation*}
    n(0) = n_{0} \quad \der{n}{r}\Big|_{r=0} = 0
\end{equation*}
allowing us to find a constant near-nucleus density $n(0)$ within a radius 
\begin{equation}
    r_{0} = \frac{3\sigma}{(4\pi G n_{0})^{1/2}}
    \label{eq:kingradius}
\end{equation}
which is known as King or core radius. At small radii ($r \lesssim r_{0}$) the density law resembles the (so-called Hubble) density law
\begin{equation*}
    n(r) = \frac{1}{\left(1+\left(\frac{r}{r_{0}}\right)^{2}\right)^{3/2}}
\end{equation*}
which fits decently the centers of many elliptical galaxies. At large radii ($r\gtrsim 15r_{0}$) the system resembles the SIS
\begin{equation*}
    n(r) \sim \left(\frac{r}{r_{0}}\right)^{-2}
\end{equation*}
which doesn't fit elliptical galaxies too well in the outer regions.
\par There is an underlying problem with all isothermal models: They have infinite total mass and size. It is easy to see why the system has infinite size:
\begin{itemize}
    \item At any given radius, stars have isotropic dispersion $\sigma$;
    \item At this radius at least some stars are moving outwards, but further out the dispersion is still $\sigma$, and stars which are moving outward the system must have infinite extent;
    \item Ultimately, this arises because $f(\epsilon) > 0$ for negative $\epsilon$, meaning the model includes unbound stars.
\end{itemize}
\par To rectify this problem, we may attempt to modify things slightly by removing the unbound stars. A viable alternative would then be the lowered isothermal (King) truncated exponential $f(\epsilon)$: It suppresses stars at large radius
\begin{equation*}
    f(\epsilon) = \left\{
        \begin{array}{ll}
            \frac{n_{1}}{(2\pi\sigma^{2})^{3/2}}\left(\exp\left(\epsilon/\sigma^{2}\right)-1\right) & \text{for}\; \epsilon > 0\\
            0 & \text{for}\; \epsilon \leq 0
        \end{array}
    \right.
\end{equation*}
where $\sigma$ is a (dispersion like) parameter. We can solve this through integration, choosing boundary conditions at $r = 0$
\begin{equation*}
    \Psi(0) = q\sigma^{2} \quad \der{\Psi}{r}(0) = 0
\end{equation*}
In the inner regions it behaves like an isothermal model, with a core King radius $r_{0}$, in the outer regions density decreases to $0$ at $r_{t}$, the tidal or truncation radius, that is the edge of the sphere for which $n(r) = 0$ for $r>r_{t}$. It is useful to define the concentration parameter $c = \log(r_{t}/r_{0})$.

\subsubsection*{Isothermal Disk}

Assume a disk that is stratified in layers parallel to its plane, so everything is a function of the vertical coordinate $z$ only. In general the velocity dispersion may change with $z$, but let's consider the case where the velocity dispersion is constant. This is called an isothermal disk. The quantity
\begin{equation*}
    E_{z} = \Phi + \frac{1}{2}v_{z}^{2}
\end{equation*}
is thus an integral of motion. The associated distribution function is 
\begin{equation*}
    f(E_{z}) = \rho_{0}(2\pi\sigma_{z}^{2})^{-1/2}\exp(-E_{z}/\sigma_{z}^{2})
\end{equation*}
which can be integrated over velocity space to obtain 
\begin{equation*}
    \rho_{z} = \rho_{0}\exp(-\Phi/\sigma_{z}^{2})
\end{equation*}
Using Poisson's equation and the standard boundary conditions we obtain the density as a function of the vertical coordinate 
\begin{equation*}
    \rho(z) = \rho_{0}\text{sech}^{2}(z/h)
\end{equation*}
where we've defined the scaleheight
\begin{equation*}
    h = \frac{\sigma_{z}}{\sqrt{2\pi G \rho}}
\end{equation*}
For $z \gg h$ we retrieve an exponentially dropping profile.
\par Is it possible in general to pass from $n$ (assuming we're able to measure it) to the distribution function $f$? No, not really. There's only a scant few of cases where that is actually possible. One is the case of spherically symmetric systems.
\par It is possible to show that given a spherically symmetric density distribution, which can be written as $\rho = \rho(\Psi)$, Eddington’s formula\footnote{It is essentially the Abel integral equation we've already encountered somewhere in our previous discussions.} yields the corresponding distribution function\footnote{This method can be extended to rotating spherical systems with $f(E, |L|)$ as well as axisymmetric systems with $f(E, L_{z})$.} $f(\epsilon)$
\begin{equation}
    f(\epsilon) = \frac{1}{\sqrt{8}\pi^{2}}\frac{\text{d}}{\text{d}\epsilon}\int_{0}^{\epsilon}\der{n}{\Psi}\frac{1}{\sqrt{\epsilon - \Psi}}\diff{\Psi}
\end{equation}

\subsection{Schwarzschild’s Orbit Superposition Method}

The method we're about to discuss is rather powerful, because no distribution function is needed for it to work.
\par The observed surface brightness is used to build a mass model for the luminous component(s) of the galaxy. Because the surface brightness corresponds to the two-dimensional projection of the underlying, 3D luminosity density, assumptions about the underlying density have to be made in order to deproject it, and values for the orientation of the galaxy with respect to the line of sight need to be chosen.
\par You have to convert the luminosity to mass, which in general is done by assuming a constant mass-to-light ratio $M/L$; this conversion then gives a mass model for the luminous component. To this you can add dark components to obtain the full gravitational potential.
\par Once you have a model for the full gravitational potential, you can use it then to create an orbit library\footnote{Template bank, is this you?} that contains the building blocks (the orbits) of a Schwarzschild model. For a given set of orbits weights, one can then predict both the surface brightness and whatever kinematic information is available.
\begin{figure}[h!]
    \centering 
    \includegraphics[width=0.8\textwidth]{img/sosm1.png}
    \includegraphics[width=0.8\textwidth]{img/sosm2.png}
    \caption{Visual representation of the SOS method. In the lower cartoon, numerical integrations of a single orbit in the adopted gravitational potential is shown. After a sufficiently long time the density of regular orbits converges to a fixes distributions.}
    \label{fgr:sosm}
\end{figure}
\par By running many instances of this basic Schwarzschild modeling procedure for different settings of the assumed parameters (inclination, mass-to-light, mass of the SBH, properties of the DM halo, etc.) and comparing the objective function for the best-fit orbital-weight distribution for each model, one can determine the overall best model.

\newpage
\section*{Some Questions to probe your Understanding}

I'll write them in Italian because I'm too lazy to properly translate them. Sooner or later I'll also add the page where you can find the answer to the respective question.

\begin{enumerate}
    \item Discutere il fenomeno della frizione dinamica: quali sono i passaggi principali per la derivazione della formula di Chandrasekhar e la sua interpretazione euristica (come si interpreta in termini di scia).
    \item Discutere il fenomeno della frizione dinamica, specificando da quali grandezze fisiche dipende. Spiegare perché un oggetto massiccio, come un ammasso globulare, tende a spiraleggiare verso il centro della nostra Galassia e stimare l’ordine di grandezza del tempo scala di decadimento orbitale.
    \item Definire il tempo di rilassamento a due corpi e discuterne il significato fisico. Valutare se le galassie e gli ammassi stellari possano essere considerati sistemi dinamicamente rilassati, confrontando i rispettivi tempi di rilassamento con le loro età caratteristiche.
    \item Rilassamento a due corpi e rilassamento violento: discutere l’origine di questi fenomeni, le differenze principali e a quali sistemi stellari si applicano.
    \item Se una galassia subisce un tidal shock, che fenomeni ti aspetti? Su che tempi scala agiscono questi fenomeni?
    \item Spiegare perché, dopo un tidal shock, una galassia tende a perdere materia più facilmente.
    \item Si discutano schematicamente i passaggi principali che portano alla formula del raggio di Hill (nel caso in cui uno dei due corpi sia molto meno massiccio dell’altro). Fare degli esempi reali (per esempio la coppia Grande Nube di Magellano-Milky Way, e la coppia Grande Nube-Piccola Nube).
    \item Nel contesto dell’interazione tra galassie si discutano i fenomeni di frizione dinamica, effetti mareali, ram pressure stripping.
    \item Equazione di Boltzmann non collisionale ed equazioni di Jeans: che cosa sono e a quali sistemi stellari si applicano. Discutere un esempio di applicazione delle equazioni di Jeans.
    \item Come possiamo stimare la massa totale nelle vicinanze del Sole tramite l’equazione di Jeans?
    \item Come possiamo stimare la massa totale di una galassia non rotante a simmetria sferica?
    \item Da cosa dipende lo schiacciamento delle galassie ellittiche? Motivare usando il teorema del viriale tensoriale. Cosa ci dice sul loro processo di formazione?
    \item Partendo da una funzione di distribuzione, discutere schematicamente come costruire una distribuzione spaziale di stelle che possa rappresentare una galassia fisica, cioè un modello realistico di una galassia esistente. Discutere inoltre se sia possibile risolvere il problema inverso, ovvero ricavare la funzione di distribuzione a partire dalla distribuzione spaziale osservata.
    \item Cosa rappresenta la funzione di distribuzione? Discutere il teorema di Jeans e mostrare come si derivano (solo i passi principali) le forme funzionali di disco e sfera isoterma.
    \item Per una galassia a simmetria sferica e statica (nell’ipotesi che la funzione di distribuzione dipenda solo dall’energia), è possibile ricavare la funzione di distribuzione a partire dal profilo di densità stellare osservato? Si discuta il metodo di Schwarzschild di sovrapposizione delle orbite.
\end{enumerate}